% TOP_PROBLEM: {"id": 2307076, "title": "Research Problems in Function Theory — Problem 7.76", "category_id": 8, "category": {"id": 8, "name": "analysis", "display_name": "Analysis", "description": "Limits, continuity, calculus, and function theory.", "slug": "analysis", "order_index": 8, "created_at": "2026-07-31T15:26:25.671Z"}, "status": "open"}
% FIRST_POSTED: null
% ENTRY_KIND: statement_only
% Imported from: batch06.tex; UnsolvedMath ID 2307076
% Compile twice: lualatex 2307076.tex
\documentclass[11pt]{article}
\usepackage{fontspec}
\setmainfont{Latin Modern Roman}
\usepackage{amsmath,amssymb,mathtools,unicode-math}
\setmathfont{Latin Modern Math}
\usepackage{xurl,hyperref}
\hypersetup{hidelinks,pdftitle={UnsolvedMath Problem 2307076}}
\newcommand{\sref}[2]{\hyperlink{source:#1:#2}{[S#2]}}
\newcommand{\cataloguescope}[1]{\paragraph{Mathematical question.}#1}
\begin{document}
% UnsolvedMath ID 2307076; source catalogue code AMR-022-7076
\setcounter{section}{2307075}
\section{Research Problems in Function Theory — Problem 7.76}\label{2307076}
{\small\sffamily UnsolvedMath ID 2307076 · Analysis}
\subsection{Definitions and mathematical statement}

\paragraph{Shared notation.}$\mathbb N=\{1,2,\ldots\}$, and $\mathbb Z,\mathbb Q,\mathbb R,\mathbb C$ denote the integers, rationals, reals and complex numbers. Any other conventions are specified locally.


For compact $E\subset\mathbb C$, continuous analytic capacity is
\[\alpha(E)=\sup\{|f'(\infty)|:f\in C(\widehat{\mathbb C}),\ f\in\mathcal O(\widehat{\mathbb C}\setminus E),\ |f|\le1,\ f(\infty)=0\}.
\]
Here $f'(\infty)=\lim_{z\to\infty}zf(z)$. On arbitrary bounded sets use inner capacity $\alpha(A)=\sup\{\alpha(L):L\subset A,\ L\text{ compact}\}$. A $C^1$ arc is the image of an injective regular $C^1$ parametrization of a compact interval; $C^{1+\varepsilon}$ additionally means its derivative is Holder continuous with positive exponent $\varepsilon$.
\paragraph{Statement recorded in the supplied source.}
If $K\subset\mathbb C$ is compact and $\alpha(K)=0$, must $\alpha(E\setminus K)=\alpha(E)$ for every compact $E\subset\mathbb C$? Retain the special case where $K$ is a $C^1$ arc. The source distinguishes this from analytic capacity and records the $C^{1+\varepsilon}$-arc case as settled. \sref{2307076}{2}
\subsection{Short English statement}
Does deleting a compact set of zero continuous analytic capacity preserve this capacity, including deletion of a continuously differentiable arc?
\subsection{Sources}
\begin{description}
\item[\hypertarget{source:2307076:1}{S1}] ULAM.AI and the UnsolvedMath Contributors, UnsolvedMath: A Curated Collection of Open Mathematics Problems, record 2307076 (AMR-022-7076). CC BY 4.0. \url{https://huggingface.co/datasets/ulamai/UnsolvedMath}.
\item[\hypertarget{source:2307076:2}{S2}] Walter K. Hayman and Edward F. Lingham, Research Problems in Function Theory (2018), Problem 7.76, its complete update and relevant preceding definitions. \url{https://arxiv.org/abs/1809.07200}.
\end{description}
\label{end:2307076}
\end{document}
